\section{Discussion and Limitations}
\label{sec:discussion}

\subsection{Evidence, not enforcement after release}
TrustTrail's strongest guarantees concern what happens at TrustTrail: what a service promised, what it asked for, what it received, and whether the record of these facts has been changed. It has no view of what a service does with the data afterwards. A service that stays inside its declaration but sells the data, or never deletes it, receives a perfect trust score. Rule~2 catches only a service that comes back after its deletion date. This is a limit that TrustTrail shares with every approach that releases plaintext attributes, including sticky policies without trusted hardware. What TrustTrail adds is that the service's promise is signed and recorded with each release, so that when misuse is discovered by other means, for example through a complaint, an audit or a leak, the institute and the student can show exactly what the service promised and received. In the terms of the GDPR and the DPDP Act, TrustTrail supports the controller in \emph{demonstrating} purpose limitation and accountability~\cite{gdpr2016,dpdp2023}; it does not prove storage limitation.

\subsection{Operator power and governance}
\label{sec:discussion-governance}
TrustTrail concentrates power in whoever operates it, and anchoring limits only part of that power. The operator controls what enters the ledger: a dishonest TrustTrail could release data without writing an entry, and anchoring cannot reveal a record that was never written. The receipt is the only counter-evidence, and it exists only if the student, or the service, keeps it. An anchor proves that an entry is unchanged since the anchor, not that it was true when written. The operator holds the receipt-signing key, the anchoring key and the identity map that links pseudonyms to students, and it can reset any service's trust; resets are recorded in the ledger, but nothing prevents them. In a university, the operator is usually the institute itself, which is also the Data Fiduciary responsible for the services it oversees, so TrustTrail does not create an independent auditor by itself.

Governance measures can narrow this gap, and none of them needs a change to the design. TrustTrail's administration can sit with an office that does not run the services, such as the institute's data protection or grievance officer, with separate people holding the administrator token and the anchoring key. Every administrative action that changes the record (a trust reset) is already a ledger entry, so it is anchored like any release. The anchor contract's address can be published, as we do in this paper, so that students, auditors and the services themselves can check that batches keep arriving and that the ledger they are shown matches the chain. Services also have an interest in watching: a service that keeps its own receipts can show that a release it received was recorded, and a service accused of a violation can check the entry against the chain.

\subsection{Linkability that remains}
Pairwise identifiers and the removal of \code{sid}, \code{session\_state} and \code{auth\_time} remove \emph{identifiers} that link a student across services. They do not remove timing. Two colluding services that compare the issue times of their tokens, or the network addresses of requests, could still guess which pseudonyms belong to the same student, especially in a small population. The anonymous profile is weaker than its name: TrustTrail releases no identifier under it, but the service's ID token from Keycloak still contains its pairwise \code{sub}, so the service can recognise a returning student. And Keycloak, as the IdP, sees every login. Privacy-preserving SSO protocols (Section~\ref{sec:related}) address IdP tracking and could be combined with TrustTrail.

The largest remaining channel is the released data itself. Department, programme and year together can single out a student in a small class, and the total profile releases name, e-mail address and roll number, which identify the student outright. Services that both hold such values can join their records without any identifier, so pairwise pseudonyms protect only students whose released attributes are common. The profile caps reduce this risk by giving low-trust services little to join on, but they do not remove it. Measuring it would need the size of the anonymity sets for each attribute combination in a real cohort, which we did not have.

\subsection{Is a blockchain needed?}
The anchor's job is to put the head of the ledger somewhere the operator cannot change it. A blockchain is one way to do this; Table~\ref{tab:anchor-alternatives} lists others. Each detects the same rewrites and truncations, provided the published value is kept by someone outside the operator's control. They differ in whom a verifier must trust, in cost, and in how much a verifier can check later.

\begin{table}[t]
\centering
\caption{Ways to fix the head of the ledger outside the operator's control}
\label{tab:anchor-alternatives}
\small
\begin{tabularx}{\linewidth}{@{}>{\raggedright\arraybackslash}p{3.4cm} >{\raggedright\arraybackslash}X >{\raggedright\arraybackslash}X >{\raggedright\arraybackslash}p{2.6cm}@{}}
\toprule
Option & Verifier must trust & Weakness & Cost \\
\midrule
Operator signs and publishes the head (web page, mailing list) & that someone archived the published value & the operator can replace the page; old values survive only if others kept them & none \\
Trusted timestamping authority (RFC~3161~\cite{rfc3161}) & the authority and its key & one more trusted party; tokens must be kept and the authority must stay available & low, often a paid service \\
Witnesses cosign the head (CT-style) & that enough witnesses are honest & needs independent organisations willing to run witnesses & operating cost for each witness \\
Public blockchain (TrustTrail) & the chain's consensus & fees; anchors are permanent and public; finality takes minutes & 254{,}408 gas per batch \\
\bottomrule
\end{tabularx}
\end{table}

We chose a public chain because, for a campus, it is the only option that needs no other organisation to cooperate. Anyone can check an anchor years later with nothing but a public RPC endpoint, as our receipt verifier does, and the contract itself enforces that batches are contiguous, so the operator cannot quietly skip or overlap ranges. The cost is a fee per batch and a dependence on the chain's continued operation. For an institute that already belongs to a federation with independent members, witness cosigning would be a reasonable alternative that avoids fees, and the ledger and receipts would not change.

\subsection{The ledger as personal data}
\label{sec:discussion-ledger}
The ledger contains no attribute values, but it is still personal data. Each entry records that a pseudonymous student received a service at a given time, and the identity map, held by the same operator, can link the pseudonym to the student. Under the GDPR, pseudonymised data that can be attributed to a person with additional information remains personal data (Recital~26)~\cite{gdpr2016}, and the DPDP Act's duties attach to it too~\cite{dpdp2023}. The current implementation keeps ledger entries indefinitely and offers no way to correct or erase them, which sits uneasily with the storage limitation that motivates the project and with a student's rights to correction and erasure (DPDP Act Section~12).

The append-only design does not have to conflict with erasure. Three measures fit it without weakening the tamper evidence. First, the institute can set a retention period for ledger entries, after which each expired entry is replaced by a tombstone that keeps only its number and hash, and its identity-map row is deleted. The chain links and the Merkle roots can still be checked, because they depend only on the hashes; what is lost is the ability to recompute an expired entry's hash from its content, which is acceptable once the entry has served its purpose. Second, erasure on request can be implemented by deleting the student's rows from the identity map: the ledger entries then carry pseudonyms that nobody can attribute, while their hashes, and therefore the chain and the anchors, are unchanged. Third, correction is handled by appending a correcting entry, never by editing an old one. We have not implemented these measures, and whether deleting the mapping is enough to make the remaining entries anonymous is a legal question that depends on what else the institute holds.

The anchored data needs separate thought. A Merkle root is a hash of hashes of entries, each of which contains a pseudonym and a timestamp. It reveals nothing about the entries without the ledger, and once the ledger entries and their mapping are deleted, the root is not linkable to anyone. Whether such hashes count as personal data while the ledger still exists is not settled; the European Data Protection Board has issued draft guidelines on processing personal data through blockchain technologies~\cite{edpb2025blockchain}. TrustTrail's design keeps the on-chain data to the minimum that anchoring needs, a root and an entry range, and keeps everything else off chain where it can be deleted.

\subsection{Deployment at an Indian institute}
In DPDP Act terms, the institute running TrustTrail is the Data Fiduciary, students are Data Principals, and services that the institute engages to process student data are Data Processors; a third-party service that decides its own purposes would be a Data Fiduciary in its own right. TrustTrail's records map onto several duties of the Act. The declaration and the dashboard give the student the processing purposes and the list of services that received their data, the information a Data Principal may request under Section~11. The signed deletion date gives the institute a recorded point at which the specified purpose ends, which it needs in order to apply the erasure duty of Section~8(7) to its processors. And the ledger and receipts are the evidence that a grievance under Sections~8(10) and~13 would need.

TrustTrail does not provide the remedy itself. A student who sees a violation on the dashboard needs a way to act on it, and the natural path is the institute's grievance redressal mechanism, which the Act requires anyway. A deployment would link each dashboard entry to that mechanism, so that a complaint carries the receipt and the entry number. Two further points matter in practice. The lawful basis for each service's processing, consent or one of the Act's legitimate uses, is a decision for the institute, and TrustTrail's declarations record the purpose but not the basis. And some students may be under eighteen, for whom the Act requires verifiable parental consent before processing; TrustTrail would need to know a student's age, or be configured to give such students only the lowest profiles until consent is recorded.

\subsection{Dependence on Keycloak internals}
The \code{sid} problem shows how a standard feature can undo a privacy property. Keycloak~26 adds \code{sid} to every token with no setting to turn it off, so we wrote a mapper against an SPI that Keycloak marks as internal. Its removal of \code{sid} breaks refresh tokens, the userinfo endpoint and logout with \code{id\_token\_hint} for service clients. The client setting that hides \code{session\_state} is deprecated in Keycloak~26.8; if a later version removes it, the value would return in the login redirect. Our browser-level check script fails in that case, so an upgrade that breaks the property would be noticed. Other IdPs would need equivalent configuration, and the same check, before TrustTrail could rely on them.

\subsection{The trust model}
The trust score is a heuristic. Its weights (15 points for a missing privacy policy, 10 for incomplete metadata, 25 per violation, one tier per 25 points) were chosen by hand and tested only with scripted clients. It ranks services by what TrustTrail can see, which is a narrow view of trustworthiness. Its strength is that it is simple and predictable: a service can see its own score through the public trust endpoint and knows that one violation costs one tier. Two design choices limit unfair effects. Requests within the service's agreement are never counted as violations after a demotion, so one mistake does not snowball. And violations expire after 90 days. Rule~3 at login depends on Keycloak's event log: if events are disabled or expire before TrustTrail reads them, the attempts are refused but not counted. TrustTrail's Keycloak service account holds the broad \code{manage-clients} role in order to change scope assignments; Keycloak's fine-grained admin permissions could narrow this.

\subsection{Key custody}
Non-repudiation rests on keys. A declaration binds a service only while its private key is not shared or stolen, and registering a key requires only the service's Keycloak client secret, so the client secret is the root of trust for key registration. TrustTrail's receipt key and the anchoring key are held by TrustTrail; whoever holds the anchoring key can add batches but cannot change or remove existing ones. TrustTrail checks declarations' attribute names, not the meaning of the purpose text: it cannot judge whether a stated purpose is honest or proportionate. That judgement belongs to the institute when it agrees a service's base tier.

\subsection{Scale and deployment}
The implementation targets a single campus. All database writes pass through one lock and one SQLite connection; the pseudonym map is rebuilt by reading at most 1{,}000 users per request; and we did not measure throughput under concurrent load. These are implementation limits rather than design limits: the ledger needs only an append-only table with a serialised append, which a server database provides, and the identity map can be updated incrementally. Keycloak runs in development mode, the services use the password grant for test scripts, and the local chain uses Hardhat's public test keys, so its anchors prove nothing to outsiders. We ran the anchoring on Sepolia to show that the design works on a public network. Sepolia is a test network, which its operators could in principle reset, and its gas is nearly free; a production deployment would use a main network, where fees apply (Section~\ref{sec:eval-anchoring}). Verifying many batches over a public RPC endpoint also takes time, about 18\,s for 15 batches in our run, so a verifier that checks often should cache batches it has already read.

\subsection{Lessons}
Three lessons from building TrustTrail may be useful to others. First, separating \emph{which} profile from \emph{what} it contains made the design simple: the IdP enforces the cap with a mechanism it already has, scope assignment, and the release service can be small. Second, every accountability mechanism we added closed a gap left by the previous one: the chain misses full rewrites, receipts cover only students who keep them, and anchors leave a window. None of them alone would have been enough, which matches the survey's observation that accountability is hard and under-served. Third, an automatic feedback loop needs care to be fair. Without the distinction between the base tier and the current tier, a single violation would demote a service, its next ordinary request would count as a new violation, and one mistake would end in blocking.
